% !TeX program = xelatex
\documentclass[UTF8,fontset=fandol,10pt,a4paper]{ctexart}
\usepackage[left=22mm,right=22mm,top=23mm,bottom=22mm,headheight=14pt]{geometry}
\usepackage{fontspec}
\setmainfont{TeX Gyre Heros}
\setsansfont{TeX Gyre Heros}
\setmonofont{Latin Modern Mono}
\renewcommand{\familydefault}{\sfdefault}
\usepackage[table]{xcolor}
\usepackage{graphicx,tikz}
\usetikzlibrary{arrows.meta,positioning,fit,backgrounds,calc}
\usepackage{booktabs,tabularx,array,enumitem,float,caption,fancyhdr,hyperref}
\definecolor{blue}{HTML}{2869A9}
\definecolor{ink}{HTML}{293340}
\definecolor{muted}{HTML}{647181}
\definecolor{line}{HTML}{C5D4E2}
\definecolor{pale}{HTML}{F0F5FA}
\definecolor{warm}{HTML}{FFF7E9}
\hypersetup{colorlinks=true,urlcolor=blue,linkcolor=blue,pdftitle={Redflag RHI 与 Vulkan 接入设计},pdfauthor={Redflag}}
\color{ink}
\setlength{\parindent}{0pt}
\setlength{\parskip}{6pt}
\linespread{1.16}
\setlist[itemize]{leftmargin=1.5em,itemsep=2pt,topsep=3pt}
\setlist[enumerate]{leftmargin=1.8em,itemsep=3pt,topsep=3pt}
\renewcommand{\arraystretch}{1.38}
\newcolumntype{Y}{>{\raggedright\arraybackslash}X}
\newcolumntype{L}[1]{>{\raggedright\arraybackslash}p{#1}}
\ctexset{section={format=\Large\bfseries\color{blue},beforeskip=16pt,afterskip=8pt},subsection={format=\large\bfseries\color{blue},beforeskip=11pt,afterskip=5pt}}
\captionsetup{font=small,labelfont=bf,textfont={color=muted},skip=7pt}
\pagestyle{fancy}\fancyhf{}
\fancyhead[L]{\small\color{muted}Redflag · RHI / Vulkan · 设计方案}
\fancyhead[R]{\small\color{muted}2026-11-21}
\fancyfoot[C]{\small\color{muted}\thepage}
\renewcommand{\headrulewidth}{0pt}
\newcommand{\note}[1]{{\small\color{muted}#1\par}}
\newcommand{\code}[1]{{\small\texttt{#1}}}
\newcommand{\ruleline}{\par\vspace{4pt}{\color{line}\hrule}\vspace{8pt}}
\newcommand{\source}[2]{{\small\href{#1}{[#2]}}}
\tikzset{
 module/.style={draw=line,fill=pale,rounded corners=2pt,align=center,inner sep=7pt,text width=3.8cm,minimum height=1.15cm,font=\small},
 smallbox/.style={module,text width=2.55cm,minimum height=.95cm,font=\footnotesize},
 wide/.style={module,text width=9.5cm},
 planned/.style={module,fill=warm,dashed},
 flow/.style={-{Stealth[length=2mm]},draw=muted,line width=.65pt},
 own/.style={flow,dashed},
 label/.style={font=\scriptsize,fill=white,inner sep=2pt,text=muted}}

% 本册的页幅、图注和代码排版。
\geometry{paperwidth=190mm,paperheight=252mm,left=21mm,right=21mm,
  top=22mm,bottom=21mm,headheight=21pt,headsep=4mm,footskip=10mm}
\usepackage{amsmath,listings,type1cm}
\definecolor{Ink}{HTML}{202E33}
\definecolor{Muted}{HTML}{4C6064}
\definecolor{Teal}{HTML}{12665F}
\definecolor{Rust}{HTML}{9E452B}
\definecolor{Rule}{HTML}{B3C6C1}
\definecolor{Soft}{HTML}{EAF3F0}
\definecolor{Paper}{HTML}{FFFFFF}
\setmainfont{TeX Gyre Pagella}
\renewcommand{\familydefault}{\rmdefault}
\linespread{1}
\setcounter{secnumdepth}{0}
\setlength{\emergencystretch}{2em}
\setlength{\parindent}{0pt}
\setlength{\parskip}{5.8pt}
\ctexset{section={format=\raggedright\sffamily\bfseries\fontsize{23}{30}\selectfont\color{Ink},
    beforeskip=0pt,afterskip=9pt},
  subsection={format=\raggedright\sffamily\bfseries\fontsize{13}{18}\selectfont\color{Teal},
    beforeskip=8pt,afterskip=2pt}}
\hypersetup{urlcolor=Teal,linkcolor=Teal,bookmarksnumbered=false,
  pdfsubject={Redflag 渲染接口、后端实现与验证}}
\fancyhf{}
\setlength{\headwidth}{\textwidth}
\fancyhead[L]{\sffamily\fontsize{9.5}{12}\selectfont\bfseries\color{Teal}Redflag\enspace{\mdseries\color{Muted}图形引擎}}
\fancyhead[R]{\sffamily\small\color{Muted}RHI / Vulkan}
\renewcommand{\headrule}{\hbox to\headwidth{\color{Rule}\leaders\hrule height .55pt\hfill}}
\fancyfoot[L]{\small\sffamily\color{Muted}\nouppercase{\leftmark}}
\renewcommand{\sectionmark}[1]{\markboth{#1}{}}
\fancyfoot[R]{\large\sffamily\bfseries\color{Teal}\thepage}
\renewcommand{\code}[1]{{\ttfamily\fontsize{10.4}{14}\selectfont #1}}
\newcommand{\readings}[1]{\par\noindent{\sffamily\fontsize{9.6}{14}\selectfont
  \color{Muted}参考\enspace #1\par}}
\newcommand{\refitem}[2]{\href{#1}{#2}}
\newcommand{\figpanel}[2]{\begin{figure}[H]\centering\input{figures/#1}\caption{#2}\end{figure}}
\DeclareCaptionFont{rhicapt}{\sffamily\fontsize{10}{14}\selectfont\color{Muted}}
\captionsetup{font=rhicapt,labelfont={bf,color=Teal},textfont=rhicapt,
  labelsep=quad,justification=raggedright,singlelinecheck=false,skip=4pt}
\setlength{\intextsep}{7pt plus 1pt minus 1pt}
\lstset{basicstyle=\ttfamily\fontsize{10.7}{15}\selectfont\color{Ink},
  backgroundcolor=\color{Soft},frame=l,rulecolor=\color{Teal},framerule=1.5pt,
  xleftmargin=5mm,xrightmargin=2mm,framexleftmargin=3mm,framesep=2pt,
  aboveskip=6pt,belowskip=6pt,columns=fullflexible,keepspaces=true,showstringspaces=false}
\newcommand{\bodytype}{\normalfont\fontsize{11.4}{17.8}\selectfont\color{Ink}}


\begin{document}
\bodytype
\section{RHI 与 Vulkan 接入设计}
Redflag 目前通过 OpenGL 把贴图传给显卡、绘制模型，并把结果显示在编辑器中。接入 Vulkan 后，场景渲染器仍负责选择可见模型和安排绘制顺序，着色器仍按相同的公式计算光照。两个后端需要共同约定的是怎样创建显卡资源、怎样把一次绘制交给显卡。这组接口通常称为渲染硬件接口（RHI）。

\subsection{资源创建与绘制命令}
材质上的底色贴图，通常在模型加载时创建一次，可以连续使用许多帧。同一模型每次出现在哪个视口、用哪份相机数据、画进哪张图，则要在绘制时决定。方案把前一件事交给 \code{RenderDevice}，把后一件事交给 \code{CommandList}。这两个名字都表示准备实现的接口。

\figpanel{interfaces}{纹理创建与模型绘制。启动时选定一个后端，两行都使用它。}

场景渲染器可以先用同一个接口在 OpenGL 中画出原来的画面。等 Vulkan 也能执行这次绘制，就可以用同一份顶点、材质和相机逐项比较结果。首版从一个窗口、一条绘制命令做起。

\newpage
\section{模型的加载与绘制}
\subsection{从资源编号到一张贴图}
场景文件用资源编号引用模型。\code{AssetManager} 根据编号找到模型，加载器读出网格与材质，材质再保存绘制所需的贴图。现有材质类中的底色贴图声明如下。下面摘出创建函数中与它对应的 OpenGL 分支。
\begin{lstlisting}
Ref<Texture2D> BaseColorMap;
\end{lstlisting}
\begin{lstlisting}
case RendererAPI::API::OpenGL:
    return CreateRef<OpenGLTexture2D>(path);
\end{lstlisting}
这里用到了\textbf{静态工厂}。加载器只要调用 \code{Texture2D::Create(path)} 就能得到一张纹理，材质可以一直用 \code{Ref<Texture2D>} 保存它。知道具体类名的是创建函数。Hazel 的公开实现也把这项选择放在 \code{Texture2D::Create} 里。

\subsection{同一个设备创建一组资源}
方案保留现有的 \code{Create} 入口，让启动时选好的设备创建配套的纹理、缓冲区和绘制管线。这是\textbf{抽象工厂}的用法。代价是每增加一种公共资源，都要决定两个后端怎样实现它。

\noindent\begin{minipage}[t]{.37\linewidth}\vspace{0pt}
以底色贴图为例，Vulkan 设备先创建图像，再按它的内存要求分配显存。材质通过图像视图采样。右边用伪代码写出创建过程，省略传入的结构体字段和像素上传。
\end{minipage}\hfill\begin{minipage}[t]{.59\linewidth}\vspace{0pt}
\begin{lstlisting}[style=besidetext]
usage = CopyDst | Sampled
spec = ColorTexture(size, SRGB, usage)
vkCreateImage(device, spec, ..., &image)
memory = AllocateFor(image)
vkBindImageMemory(device, image, memory, 0)
view = CreateColorView(device, image)
\end{lstlisting}
\end{minipage}

\subsection{主画面与阴影的可见模型}
\code{SceneRenderer} 从相机挑选主画面里的模型，还要从光源再挑一次投影物。相机外的物体仍可能在画面里投下阴影。这项选择留在场景渲染器里，两个后端只接收挑好的网格、材质和目标图像。
\readings{\refitem{https://github.com/TheCherno/Hazel/blob/1feb70572fa87fa1c4ba784a2cfeada5b4a500db/Hazel/src/Hazel/Renderer/Texture.cpp}{Hazel 的 Texture.cpp}\quad
\refitem{https://www.informit.com/articles/article.aspx?p=1398599}{《设计模式》抽象工厂}\quad
\refitem{https://docs.vulkan.org/refpages/latest/refpages/source/vkCreateImage.html}{vkCreateImage}\quad
\refitem{https://cdn.boomingtech.com/games104_static/upload/GAMES104_lecture02Layered\%20Architecture\%20of\%20Game\%20Engine.pdf\#page=26}{GAMES104 第 2 讲}}

\newpage
\section{场景颜色与后处理}
场景渲染器只需调用一次 \code{PostProcess::Render}。后处理在内部提取亮部、逐级缩小和放大，再把光晕加回场景。将这些步骤藏在一个入口后面，就是\textbf{外观模式}的一种用法。各级泛光图由后处理对象持有。

现有调用要传入 OpenGL 纹理编号、宽和高。拟议接口改为传入视口持有的两张图，一张用来读取场景颜色，另一张用来写合成结果，尺寸可以直接从图像取得。
\begin{lstlisting}
postProcess.Render(commands, view.SceneColor(), view.Output());
editor.DrawImage(view.Output());
\end{lstlisting}
编辑器调用 \code{DrawImage} 显示结果。OpenGL 界面实现取出纹理编号，Vulkan 界面实现向 ImGui 登记图像视图和采样器。

\subsection{面板放大后，旧画面何时释放}
假设场景面板原来是 $640\times360$ 像素。上一帧已交给显卡，里面有一条界面绘制命令，要读取这张小图，画出面板里的方块。此时把面板拖大，下一帧就要准备一张 $960\times540$ 的图。
\figpanel{texture-lifetime}{同一台相机看到的方块。两个矩形表示显存中两张不同尺寸的场景图，按尺寸比例示意。}
新图建好就能安排下一帧。上一帧的命令里写的仍是小图，等显卡读完它，才能注销小图在 ImGui 里的记录，释放它占用的显存。

\newpage
\section{合成画面的绘制命令}
现有的 \code{PostProcess::Composite} 绑定场景颜色和泛光图，用一个覆盖整幅画面的三角形运行合成着色器。迁移后，同一次合成可以写成下面几行。这里的 \code{inputs} 包含两张输入图和曝光值。
\begin{lstlisting}
commands.BeginRendering(output);
commands.BindPipeline(m_CompositePipeline);
commands.BindVertexBuffer(m_FullscreenVertices);
commands.BindResources(inputs);
commands.Draw(3);
commands.EndRendering();
\end{lstlisting}

\subsection{绘制目标、管线与输入}
\code{output} 指明写入哪张图、哪一层，以及原来的颜色是否要保留。最后合成会覆盖所有像素，可以舍弃原来的颜色。把较小一级的泛光加回上一层时，却要保留上一层已有的亮度。如果接口漏掉这个区别，画面上可能只剩最后加进去的光。

\code{Pipeline} 保存可重复使用的规则，包括着色器、目标格式、深度测试和混合方式。\code{BindResources} 选择这次实际读取的数据。两个视口可以共用一条合成管线，同时保留各自的颜色图和曝光值。只有需要挑选纹理的某一级或立方体的某一面时，接口才额外传入图像视图。

\subsection{OpenGL 调用与 Vulkan 录制}
当前 \code{RenderCommand::DrawIndexed} 直接转交给 OpenGL。Vulkan 则把目标、管线、输入和顶点数记入命令缓冲区 \code{cmd}，稍后一起提交。经典的\textbf{命令模式}会把每次请求做成独立对象，这里的命令列表直接使用后端提供的记录能力。

\noindent\begin{minipage}[t]{.37\linewidth}\vspace{0pt}
右边是上面六行在 Vulkan 端的伪代码。先让输出图可写，再画三角形，最后让界面能采样写好的颜色。两次转换都由后端记录图像屏障；输入贴图已准备好，管线的颜色格式与输出一致。
\end{minipage}\hfill\begin{minipage}[t]{.59\linewidth}\vspace{0pt}
\begin{lstlisting}[style=besidetext]
ToColorAttachment(cmd, output)
target = ColorTarget(output, Discard, Store)
vkCmdBeginRendering(cmd, target)
BindComposite(cmd, pipeline, vertices, inputs)
vkCmdDraw(cmd, 3, 1, 0, 0)
vkCmdEndRendering(cmd)
ToShaderRead(cmd, output)
\end{lstlisting}
\end{minipage}
\readings{\refitem{https://docs.vulkan.org/samples/latest/samples/extensions/dynamic_rendering/README.html}{Khronos 动态渲染示例}\quad
\refitem{https://gameprogrammingpatterns.com/command.html}{《Game Programming Patterns》命令模式}}

\newpage
\section{窗口与第一个三角形}
现在 \code{WindowsWindow::Init} 创建 OpenGL 上下文，\code{RenderCommand} 的实现也在启动时固定为 OpenGL。迁移的第一步是在创建窗口前选一次后端，默认沿用 OpenGL。选 Vulkan 时，让 GLFW 创建不带 OpenGL 上下文的窗口，按它提供的扩展列表创建 Vulkan 实例，再为这个窗口创建显示表面。

\subsection{本机设备与工程设置}
项目已有的设备查询记录显示，RTX 5070 Ti Laptop GPU 报告 Vulkan 1.4.351，并提供能同时绘制和呈现的队列。首版按 Vulkan 1.3、C++17 和 GLSL 开发。创建设备时查询并启用三项功能。动态渲染让程序在开始绘制时指定目标图像。同步 2 用来安排前一次写入和后一次读取。时间线信号量用一个递增数值表示显卡执行到了哪一次提交。

前两项通过 Vulkan 1.3 的 \code{dynamicRendering} 和 \code{synchronization2} 启用，后一项通过 1.2 的 \code{timelineSemaphore} 启用，窗口还需要交换链扩展。本机记录的 SDK 为 1.3.275.0。若生成工程的终端尚未配置它，可按以下方式指定位置。
\begin{lstlisting}
$env:VULKAN_SDK = 'C:\VulkanSDK\1.3.275.0'
.\GenerateProjects.bat
\end{lstlisting}
Premake 还需要为 Vulkan 后端链接 \code{vulkan-1.lib}。工程已有的 shaderc 用于编译着色器。

\subsection{首帧的验证}
测试程序和 Forge 共用同一套开窗与设备创建代码。先画一个三个角颜色不同的三角形，在 RenderDoc 中找到对应的绘制，检查顶点和窗口图像。再测试拖动、最小化、恢复和退出。窗口按 GLFW 报告的实际像素大小取图，宽或高为零时继续处理事件，恢复后再绘制。OpenGL 在帧末交换缓冲区，Vulkan 在帧末提交窗口图像供显示。首帧通过仍需这项运行测试。
\readings{\refitem{https://docs.vulkan.org/tutorial/latest/03_Drawing_a_triangle/00_Setup/01_Instance.html}{Khronos Vulkan Tutorial · Instance}}

\newpage
\section{模型矩阵与材质贴图}
画出三角形后，把同一份网格画两次，两次使用不同的位置、颜色和实体编号。CPU 可能在 GPU 开始绘制第一个模型之前，就已准备好第二个模型。如果两次都写进同一处内存，显卡最后可能把两个模型画在一起。

\subsection{每次绘制保留一份矩阵}
首版轮流使用两份帧记录，每份保存命令、待上传数据和用到的资源。每画一个模型，就从当前帧的上传区取一段新空间并复制矩阵。假设设备要求起始位置按 256 字节对齐，矩阵占 64 字节，第一份放在偏移 0，第二份便从 256 开始。实际间隔要查询设备。只有显卡的完成值达到本次提交编号，程序才会复用这份空间。已经录制但尚未提交的命令也持有资源，若放弃录制便释放这些引用。

\subsection{灯光数据与材质贴图}
模型片元着色器的灯光和底色图都用到了数字 3。下面摘出资源声明，省略了灯光块的成员。OpenGL 给数据块和纹理分别编号。Vulkan 按组查找资源，直接沿用这两处声明会把它们放到同一组的同一项。
\begin{lstlisting}
layout(std140, binding = 3) uniform Lights { /* ... */ };
layout(binding = 3) uniform sampler2D u_BaseColorMap;
\end{lstlisting}
方案将灯光归入各模型共用的帧数据，底色图归入材质，每次变化的矩阵和编号归入对象。编译后要读取资源类型、所在组和位置，对照 CPU 与各着色阶段，再实际画出两个模型。即使单个阶段通过编译，材质也可能读错图。

\subsection{已有的着色器检查}
此前将 15 个文件拆成 30 个阶段，以默认的非实例化版本运行 \code{glslc} 和 \code{spirv-val}。17 个通过，13 个在编译时被拒绝。其中 10 个缺少跨阶段输入输出的位置，\code{FlatColor} 的两个阶段使用了 Vulkan 不接受的独立数值参数，2D 精灵片元阶段缺少贴图数组的位置。迁移先覆盖测试模型，再验证实例化版本。两个后端分别缓存编译结果，区分源码及其包含文件、编译选项和编译器版本。

\newpage
\section{深度换算与鼠标拾取}
现有相机矩阵按 OpenGL 约定计算。顶点的裁剪坐标是 $(x,y,z,w)$，透视除法后，深度 $z/w$ 从近裁面的 $-1$ 变化到远裁面的 $1$。Vulkan 需要同两个端点落在 $0$ 和 $1$。要让左端移到零并保持右端不动，只需把原来的数加一，再除以二。
\figpanel{clip-depth}{上下两行是同一顶点在两套约定下的数值。中间刻度表示数值区间的中点，透视相机下它对应的视距取决于投影。}
顶点着色器输出时还没除以 $w$。把下式两边乘以 $w$，便得到可以直接写入顶点位置的那一行代码。
\[
  \frac{z_{\mathrm{vk}}}{w}=\frac{1}{2}\left(\frac{z_{\mathrm{gl}}}{w}+1\right)
  \qquad\Longrightarrow\qquad z_{\mathrm{vk}}=\frac{z_{\mathrm{gl}}+w}{2}.
\]
\begin{lstlisting}
gl_Position.z = 0.5 * (gl_Position.z + gl_Position.w);
\end{lstlisting}
只在 Vulkan 最终输出顶点时换算一次。CPU 的剔除和光源投影仍使用原矩阵。在默认深度范围下，从深度图读回 $d$ 后，用 $2d-1$ 配合原逆投影矩阵，可还原 SSAO 与 RSM 所需的位置。

\subsection{像素从哪一行开始}
Vulkan 视口取 $y=H$、高度 $-H$，其中 $H$ 是图像的像素高度。用原投影得到的纵坐标 $v$ 采样屏幕图像时换成 $1-v$，模型自带的贴图坐标沿用原值。拾取先将鼠标换到视口像素。顶行向下数的第 $y$ 行，在 OpenGL 左下原点的坐标下为 $H-1-y$。项目的 $32\times32$ 测试在底行起算的 $(8,16)$ 读到实体 71，在顶行起算的同一位置应读 $(8,15)$。第二个点 $(24,15)$ 在实例化绘制时为 72，单体绘制时为背景 $-1$。
\readings{\refitem{https://docs.vulkan.org/guide/latest/depth.html}{Khronos · Depth}\quad
\refitem{https://www.realtimerendering.com/Real-Time_Rendering_4th-Appendices.pdf\#page=19}{RTR4 · 附录 A.4}}

\newpage
\section{编辑器视口与窗口显示}
场景先画进一张可以保存高亮度值的彩图。后处理提取亮部做泛光，再调曝光，把结果写进视口的最终颜色图。ImGui 读取这张图，将它和工具栏、其他面板一起画进窗口。窗口用来轮流显示的图像称为交换链图像。拖动场景面板只影响视口尺寸，只有窗口本身发生变化才需要考虑重建交换链。

\subsection{显卡画完与窗口显示}
当前随库的 ImGui Vulkan 实现接受 \code{VkRenderPass}。计划将它留在界面代码内，场景绘制使用动态渲染。先接主窗口中的停靠面板，统一登记视口图和工具栏图标。若表面格式或最小图像数改变，重新配置界面后端与相关登记。可拖出主窗口的独立窗口以后各自管理交换链。

显卡完成一次绘制后，工作帧就能复用其命令空间，窗口系统却可能仍在显示那张图。获取新图像所用的信号随工作帧保存，供呈现等待的信号随交换链图像保存。重新取得同一张图并等到获取完成，才可复用其呈现信号。若取得图像时报告交换链过期，按新尺寸重建，未提交的工作不记完成值。旧交换链的回收还要确认呈现完成，可用呈现完成扩展或 Khronos 给出的延迟回收办法。
\readings{\refitem{https://docs.vulkan.org/guide/latest/swapchain_semaphore_reuse.html}{Khronos · Swapchain Semaphore Reuse}}

\subsection{连续两次读写的顺序}
阴影阶段先写深度图，随后模型着色器采样它。每个阶段写明读哪张图、写哪张图及其哪一层，后端才知道何时等写入完成、何时切换到可采样的布局。实体编号图只需复制鼠标所在的一个像素，等复制完成再读整数。鼠标在视口外或图像已经更换，便忽略这次读回。RSM 的两张历史图轮流读写，前一次读取结束后才能覆盖，相机切换或尺寸变化时清除旧历史。

目前合成着色器已经为显示编码一次。首版优先使用 UNORM 窗口格式与 sRGB 显示色彩空间。若窗口只支持自动编码的 sRGB 图像，先在 UNORM 图中合成界面，再解码并写入最终目标，让硬件完成最后一次编码。检查相同曝光下的灰块和工具栏，可发现重复编码造成的发白。
\readings{\refitem{https://docs.vulkan.org/guide/latest/synchronization_examples.html}{Khronos 同步示例}\quad
\refitem{https://learnopengl.com/Advanced-OpenGL/Framebuffers}{LearnOpenGL · Framebuffers}}

\newpage
\section{一帧的提交与资源回收}
首版计划把场景与主窗口界面一起交给同一条图形队列。下面取成功获取图像的一帧说明。伪代码保留主要 Vulkan 函数名，省略结构体字段。

\subsection{取得窗口图像}
工作帧先等它的上一次提交结束，再复用命令空间。\code{vkAcquireNextImageKHR} 取得的序号 \code{i} 决定界面画进哪张窗口图。录制时工作帧持有视口输出图和界面登记。界面画完后，再把窗口图切换到可呈现的布局。
\begin{lstlisting}
WaitForPreviousUse(frame)
i = AcquireWindowImage(swapchain, frame.imageAvailable)
BeginCommands(frame.cmd)
RecordSceneAndPostProcess(frame, viewport)
RecordEditor(frame, viewport.output, swapchain.images[i])
EndCommands(frame.cmd)
\end{lstlisting}

\subsection{绘制队列与呈现队列}
本次提交等待取图成功，画完后发出两个信号。二值信号通知窗口可以显示第 \code{i} 张图，时间线信号让 CPU 知道编号 \code{serial} 对应的整次提交已完成。示例取保守的等待和发信范围，涵盖全部命令。
\begin{lstlisting}
serial = NextSubmissionNumber()
work = Submit(frame.cmd, wait = frame.imageAvailable,
    signal = [presentReady[i], timeline.At(serial)])
vkQueueSubmit2(graphicsQueue, 1, &work, nullFence)
frame.lastSubmit = serial
vkQueuePresentKHR(presentQueue, Present(i, presentReady[i]))
\end{lstlisting}

\subsection{释放上一帧用过的资源}
只有提交成功才更新编号。CPU 用 \code{vkGetSemaphoreCounterValue} 查询完成值。达到 \code{frame.lastSubmit} 后，工作帧放下它持有的引用。若换下的视口图已没有其他使用者，就可以注销界面登记、销毁视图和图像，再归还显存。窗口图像的回收还要等呈现系统用完。
\readings{\refitem{https://docs.vulkan.org/refpages/latest/refpages/source/vkQueueSubmit2.html}{Khronos · vkQueueSubmit2}\quad
\refitem{https://docs.vulkan.org/samples/latest/samples/extensions/timeline_semaphore/README.html}{Timeline semaphore}}

\newpage
\section{迁移顺序与验收场景}
\subsection{窗口与模型}
公共声明计划放在 \code{Renderer/RHI}，具体实现放在 \code{Platform/OpenGL} 和 \code{Platform/Vulkan}。第一步合并启动时的后端选择，确认 Forge 能用 OpenGL 正常打开，再让同一个窗口入口在两端各画一个三角形，保存 RenderDoc 帧。已有设备查询和离线着色器检查，Vulkan 首帧仍待实施。

第二步接入 \code{Buffer}、\code{Texture2D} 和 \code{Shader} 的创建与绘制。先画无光照的底色贴图，再画两个材质、位置不同的模型。复用已有深度预写与拾取测试，测四角、近裁面、启用背面剔除后的正面，以及被透明贴图裁掉的叶片。两个后端导出同格式的像素数据。

\subsection{编辑器与渲染效果}
第三步让视口持有尺寸和输出，面板直接接收图像对象。开两个曝光不同的视口，连续拖动尺寸、点击上边缘，在提交后删除刚绘制的模型并退出。开启 Vulkan 验证层和同步验证，检查最后一次提交何时用完旧纹理。

第四步依次接回深度预写、阴影、环境遮蔽、环境光预计算、RSM、泛光以及 2D 的精灵、圆、线和粒子。Vulkan 首版拟用 16 个 2D 贴图位置，包含一张白图。填齐数组后，越过容量便拆批，检查下一批的贴图与实体编号。OpenGL 保留原批容量，是否扩大 Vulkan 容量待测。

\subsection{用同一帧比较结果}
Cornell Box 固定相机、灯光和曝光，分辨率取 $1024\times768$，在第 12 帧检查相同位置的编号、深度和线性颜色。RSM 的采样方向依赖像素中心，两端先统一坐标，再比较同一块墙的反光。红墙用例沿用 $\Delta R>1.5\Delta G$，绿墙用例沿用 $\Delta G>\Delta R$，差值取同一白表面开关间接光前后的变化。Sponza 另测多材质与透明裁剪。

\path{tests/run_depth_tests.cmd} 测深度和拾取，\path{tests/run_scene_tests.cmd} 测整帧，\path{tests/run_rsm_scene_tests.ps1} 测彩色反光。保留版本、配置、驱动、读回值与捕获帧。性能在同一配置下另行预热和采样。

\newpage
\section{设计所依据的资料}
\subsection{Hazel 与 Piccolo 的源码}
在 Hazel 的 \code{Texture.cpp} 中，两个 \code{Texture2D::Create} 分别接受规格与路径，内部按当前 API 创建具体纹理。\code{RenderCommand.h} 则展示了如何把绘制转交给已有后端。Redflag 沿用了这种简洁入口。对于 Vulkan 录制命令和保留在途资源，这两处公开代码还需要配合其他资料。
\readings{\refitem{https://github.com/TheCherno/Hazel/blob/1feb70572fa87fa1c4ba784a2cfeada5b4a500db/Hazel/src/Hazel/Renderer/Texture.cpp}{Hazel · Texture.cpp}\quad
\refitem{https://github.com/TheCherno/Hazel/blob/1feb70572fa87fa1c4ba784a2cfeada5b4a500db/Hazel/src/Hazel/Renderer/RenderCommand.h}{RenderCommand.h}}

Piccolo 的 \code{RenderSystem::initialize} 组装 Vulkan RHI 和渲染资源，资源对象负责上传环境光等数据。它的 RHI 还提供创建缓冲区、分配描述符集合等细粒度调用。对 Redflag 有用的是这处分工。模型与场景决定需要什么，设备处理怎样在显卡上创建它。首版只公开合成和模型绘制实际用到的操作。
\readings{\refitem{https://github.com/BoomingTech/Piccolo/blob/f5053707fed4d3f94d270a436fb0d3a8ae54e3e5/engine/source/runtime/function/render/render_system.cpp}{Piccolo · RenderSystem}\quad
\refitem{https://github.com/BoomingTech/Piccolo/blob/f5053707fed4d3f94d270a436fb0d3a8ae54e3e5/engine/source/runtime/function/render/interface/rhi.h}{RHI 声明}}

\subsection{渲染流程与图像读写}
Khronos 的教程把场景剔除、安排绘制阶段、生成命令与设备执行分开。Unity 6 的 URP 示例在一个阶段分别声明输入纹理和输出图像，再提供实际绘制函数。Redflag 可以先将固定顺序中的阴影、主绘制、后处理各自读写哪张图写清楚。这样就能处理图像之间的依赖，等确有动态排程的需求时再考虑完整的渲染图。
\readings{\refitem{https://docs.vulkan.org/tutorial/latest/Building_a_Simple_Engine/Engine_Architecture/05_rendering_pipeline.html}{Khronos · Rendering Pipeline}\quad
\refitem{https://docs.unity.com/en-us/engine/6000.0/manual/render-pipelines/universal-render-pipeline/customizing-urp/render-graph/write-render-pass}{Unity URP · 自定义渲染阶段}\quad
\refitem{https://docs.vulkan.org/tutorial/latest/_images/images/rendering_pipeline_flowchart.png}{教程原图}}

\subsection{场景对象与渲染器的分工}
Gregory 在《游戏引擎架构》首版第 14.6.2 节用网格实例说明，对象只更新自己的变换和可见性，渲染系统可以收集这些实例，等对象更新结束后统一绘制。Redflag 将模型选择保留在场景渲染器中，就可以同时考虑相机和光源。第四版另核对了卷二的目录，第 11、12 章分别讨论渲染、光照与后处理。这里引用的网格例子出自首版原文。
\readings{\refitem{https://www.gamedeveloper.com/programming/book-excerpt-game-engine-architecture}{Gregory · 首版第 14 章书摘}\quad
\refitem{https://www.routledge.com/Game-Engine-Architecture-Volume-II-Graphics-Motion-and-Sound/Gregory/p/book/9781041162575}{第四版卷二目录}}

GAMES104 第 2 讲第 12–14 页讲资源导入与寿命，第 26、30 页讲 RHI 与上层的依赖。GAMES101 的投影和 RTR4 的齐次坐标附录用于核对深度公式。其他教材、课件以及具体核对的代码文件，见项目的图形学资料索引。
\end{document}
